% !TeX root = ../../Book3.tex
\section{椭圆曲线的解析表示}
\label{b3:sec:2204}

微分方程把一对亚纯函数放在一条代数曲线上。但“满足方程”尚不等于得到曲面的参数化，还需要验证无穷远点、局部坐标、单射和满射。

\subsection{三次方程}
\label{b3:subsec:220401}

把仿射方程 \(y^2=4x^3-g_2x-g_3\) 齐次化为
\[
E_\Lambda:\quad Y^2Z=4X^3-g_2XZ^2-g_3Z^3.
\]
在 \(Z=0\) 上只有一点 \(O=[0:1:0]\)。有限点的奇点须同时满足 \(y=0\)、\(12x^2-g_2=0\) 及曲线方程，等价于右侧三次多项式有重根。前节已经排除重根；在 \(O\) 处，对 \(Z\) 的偏导为一，也非奇点。因此这是一个光滑紧复曲线。

\subsection{判别式}
\label{b3:subsec:220402}

多项式 \(4X^3-g_2X-g_3\) 的无重根条件等价于
\(D=g_2^3-27g_3^2\ne0\)。这里把 \(D\) 称为该 Weierstrass 方程的\term[tuoyuan-quxian-panbieshi]{判别式}[discriminant]；有的代数规范另带固定常数，使用具体公式时应核对归一化。

\ProseParagraphBreak

换格为 \(a\Lambda\) 时，由绝对收敛格和可得
\[
\wp(az;a\Lambda)=a^{-2}\wp(z;\Lambda),\quad
g_2(a\Lambda)=a^{-4}g_2(\Lambda),\quad g_3(a\Lambda)=a^{-6}g_3(\Lambda).
\]
因此判别式的缩放次数为十二。这个整数将成为下一章模形式权十二的来源，并不是任意选择的指数。

\subsection{\texorpdfstring{\(\mathbb C/\Lambda\)}{C/\AnalysisPdfCapitalLambda}}
\label{b3:subsec:220403}

商空间 \(\mathbb C/\Lambda\) 由平行四边形粘合对边得到。每个小于半个最短周期的圆盘都与其非平凡格平移不相交，其投影给出复坐标。坐标变换为平移，因而全纯；闭平行四边形覆盖商空间，故它紧致。

\ProseParagraphBreak

对 \(z\notin\Lambda\) 定义 \(\Phi(z)=[\wp(z):\wp'(z):1]\)，对格点定义 \(\Phi(z)=O\)。在原点附近除以 \(\wp'\)，坐标为
\([\wp/\wp':1:1/\wp']=[-z/2+O(z^5):1:-z^3/2+O(z^7)]\)，所以该赋值全纯，且在原点为浸入。它按周期下降为商曲面上的全纯映射。

\subsection{解析统一化}
\label{b3:subsec:220404}

\begin{theorem}\label{b3:thm:22-uniformization}
映射 \(\Phi:\mathbb C/\Lambda\to E_\Lambda\) 是双全纯映射。
\end{theorem}
\begin{proof}
若两个有限点具有同一 \(x\) 坐标，则对应参数为 \(z\) 与 \(-z\)。它们的 \(y\) 坐标相反；若还相等，则 \(\wp'(z)=0\)，于是 \(z\) 为半周期，\(z\equiv-z\pmod\Lambda\)。无穷远点只来自格点，故单射。对任意 \((x,y)\in E_\Lambda\)，次数二保证 \(\wp\) 取到 \(x\)；微分方程给导数值为两个平方根之一，换 \(z\) 为 \(-z\) 即得到指定 \(y\)，所以满射。
在 \(\wp'\ne0\) 处可用 \(x\) 作局部坐标；在半周期处，\(\wp'\) 的零点为单零点，故 \(y\) 是局部坐标；在 \(O\) 处用上面 \(\wp/\wp'\) 的非零一阶项。逆函数定理在三种情形都给全纯局部逆，故全局逆也全纯。
\end{proof}
这里证明的是由格构造的曲线的解析表示。反过来从任意光滑复三次曲线恢复周期格，需要积分全纯微分并核对周期；这一一般方向可结合\appref{b3:app:A}的 Riemann 曲面理论继续学习，不能仅由上述显式代入推出。

\subsection{群律的解释}
\label{b3:subsec:220405}

通过 \(\Phi\) 将商群加法搬到曲线上，\(O\) 为单位元，\((x,y)\) 的逆为 \((x,-y)\)。其几何意义可由除数关系读出：一般直线与三次曲线交于三点，按重数记为 \(P,Q,R\)，相应线性函数在环面上只有原点三阶极点，因此三个参数之和模格为零。故 \(P+Q\) 是第三交点关于 \(x\) 轴的对称点。

\ProseParagraphBreak

当两点的横坐标不同时，直线斜率为 \(m=(y_2-y_1)/(x_2-x_1)\)，代入三次方程并比较根的和得
\[
x(P+Q)=m^2/4-x_1-x_2.
\]
对应纵坐标取第三交点的负值。两点合并时斜率改为 \((12x_1^2-g_2)/(2y_1)\)；纵坐标为零的切线情形给 \(2P=O\)。这些公式来自已经存在的群结构，结合律由参数相加自动保证，无需再对繁长有理式作三重展开。

\ProseParagraphBreak

还可确定整个椭圆函数域。若 \(f\) 为偶椭圆函数，它在 \(\wp\) 的每条纤维上取同一值，因此下降为球面上的亚纯函数，必为 \(\wp\) 的有理函数。在半周期处，局部参数变换 \(u\mapsto-u\) 使偶 Laurent 展开只含 \(u^{2k}\)，保证下降函数确实亚纯。对一般 \(f\)，偶部 \((f(z)+f(-z))/2\) 为 \(R(\wp)\)，奇部除以 \(\wp'\) 后也为偶函数，故
\[
f(z)=R(\wp(z))+\wp'(z)S(\wp(z))
\]
其中 \(R,S\) 有理。这说明曲线方程已经包含全部椭圆函数，而不只是两个偶然选择的函数。

\ProseParagraphBreak

三次曲线上的微分也可以通过参数直接理解。在非分歧点，\(x=\wp(z)\)、\(y=\wp'(z)\)，所以
\[
 \frac{\mathrm{d}x}{y}=\mathrm{d}z.
\]
在半周期处，两边仍全纯：若 \(z=z_0+t\)，则 \(x=e_j+\tfrac12\wp''(z_0)t^2+O(t^3)\)、\(y=\wp''(z_0)t+O(t^2)\)，分子与分母的一阶因子相消，商趋于 \(\mathrm{d}t\)。在无穷远，\(x=z^{-2}+O(z^2)\)、\(y=-2z^{-3}+O(z)\)，又得 \(\mathrm{d}x/y=(1+O(z^4))\mathrm{d}z\)。因此这个看似含分母的微分在整个紧曲线上无极点、无零点。
\ProseParagraphBreak
环面上的任意全纯微分提升到平面后写成 \(h(z)\mathrm{d}z\)。对格平移不变要求 \(h\) 是整椭圆函数，故常数。因此全纯微分空间恰为一维，由 \(\mathrm{d}z\) 生成。沿两条基本闭曲线积分得到 \(\omega_1,\omega_2\)，周期格遂从微分中恢复出来。这个结论是在已知格环面上直接计算；一般紧 Riemann 曲面的微分空间维数需要附录 A 的进一步理论。
\ProseParagraphBreak
二重映射 \(\wp\) 的四个分歧点也由局部式看得清楚。三个非零半周期在有限值 \(e_1,e_2,e_3\) 上有局部参数平方；原点的二阶极点在目标无穷远坐标 \(1/\wp\) 下同样是平方。其他点导数非零，局部为双全纯。因而三次曲线写成 \(y^2=4(x-e_1)(x-e_2)(x-e_3)\) 时，选择一个 \(x\) 通常有两种符号的 \(y\)，绕一个单根则交换这两种选择。代数方程的平方根分支与环面上的半周期是同一现象的两种坐标描述。
\ProseParagraphBreak
矩形格能进一步给出实图像。设 \(\omega_1>0\)、\(\omega_2\) 纯虚且虚部正，格关于共轭不变，由级数有 \(\overline{\wp(z)}=\wp(\bar z)\)。实轴上 \(\wp\) 实值，从原点附近的正无穷开始下降；在 \((0,\omega_1/2)\) 上导数没有零点，所以直到半周期才达到第一个临界值。它必须是三个实根中的最大者，否则下降途中先遇到别的根就会先有导数零点。沿虚轴，\(\wp(\mathrm{i}t)\sim-t^{-2}\)，从负无穷上升，首次临界值同理是最小根。第三个半周期给中间根。三个值确为实数，因为共轭半周期与原半周期相差格点。由此不用数值画图，就能确定矩形格的三个分歧值顺序及实椭圆曲线的两支。
\ProseParagraphBreak
上述顺序依赖矩形格假设。对一般倾斜格，三个分歧值可以全为非实复数，谈“最大根”便没有意义；但无重根、四个二阶分歧点和全纯微分的结论仍成立。几何图像有助于选择积分路径，却不能替代复参数情况下的局部代数检验。
